% Living draft. Numbers come from research/journal/LOG.md,
% research/logs/analysis/BASELINE_REVIEW.md,
% research/logs/analysis/CYCLE1_REVIEW.md,
% research/logs/analysis/CYCLE2_REVIEW.md, and
% research/logs/analysis/CYCLE3_REVIEW.md.
% This is cycle 3 of an unfinished study. The goal checks have not passed.
\documentclass[11pt]{article}

\usepackage[T1]{fontenc}
\usepackage[utf8]{inputenc}
\usepackage{lmodern}
\usepackage{microtype}
\usepackage{graphicx}
\usepackage{booktabs}
\usepackage{amsmath}
\usepackage[margin=1in]{geometry}
\usepackage[backend=biber,style=authoryear]{biblatex}
\addbibresource{references.bib}

\title{Teacher-Student Learning of a Predictive Locomotion Policy from a Reactive Controller}
\author{WaLTER predictive-policy project}
\date{Draft, 14 October 2026}

\begin{document}
\maketitle

\begin{abstract}
The full-body PD controller on WaLTER~FS is a reactive teacher: it tumbles only after a speed-lag integral crosses a gate, and it does not read \texttt{terrain\_scan}.
The student observes a torso height scan.
Cycle~3 changed the predictive knee term so a forward command and an open scan gate still pay when measured speed is at or below zero, and so a downward step opens the same \(0.04\,\mathrm{m}\) to \(0.08\,\mathrm{m}\) gate as an upward step.
Training then continued from the cycle-2 checkpoint for 500 further iterations at seed 42 with 2048 environments.
The resulting policy keeps flat knee rates near zero (means \(0.016\) and \(0.014\,\mathrm{rad/s}\)) and flat speed near the \(1\,\mathrm{m/s}\) command (mean \(1.027\,\mathrm{m/s}\) on both difficulties).
That is about \(0.33\,\mathrm{m/s}\) faster than the PD cruise of \(0.697\,\mathrm{m/s}\), which is the PD undershoot and is not scored as a flat failure.
Hard pyramid stairs are cleared by pitching.
Hard inverted stairs and the hard box grid both start the knees earlier than the PD, and neither cell is cleared.
A scan ablation on hard pyramid stairs moves knee-rate onset from \(0.680\,\mathrm{m}\) to never.
One seed was trained, so two-seed agreement is not met.
This is cycle~3 of an unfinished study.
The goal checks have not passed.
\end{abstract}

\section{Introduction}
\label{sec:introduction}

WaLTER~FS is trained by pairing a full-body PD controller with a student policy.
The PD is a reactive teacher.
It pitches, rolls, drives, and tumbles from body state and the velocity command.
Tumbling starts after a speed-lag integral crosses a threshold.
The PD never reads \texttt{terrain\_scan}.

The student observation includes a torso height scan the PD does not read.
A student can raise reward by cloning the PD and ignoring that scan.
The result this paper is set up to test is a policy that matches the PD on flat ground---upright, wheel-driven, knees quiet---and that tumbles earlier than the PD when the forward scan shows stairs or broken ground.

Cycle~0 recorded the PD on one cell at a time.
Those clips are the reference, because the previous teacher checkpoint does not load into the current 16-action, 383-observation environment.
Cycle~1 is the first training run of the gated objective.
Cycle~2 left that gate in place and continued the same seed.
Cycle~3 changed the stall scale and counted drops in the step measure, then continued again from the cycle-2 checkpoint.
None of the cycles passes the goal checks.
Two seeds have not been compared.
This draft reports that unfinished study.
It does not claim the result.

\section{Related Work}
\label{sec:related}

The usual teacher--student split gives the teacher more information than the student.
\textcite{lee2020quadrupedal} train a privileged teacher in simulation and imitate it with a blind student.
\textcite{chen2020learning} use the same pattern for driving: a privileged agent sees the layout, and a vision-only student copies it.
\textcite{miki2022perceptive} keep a privileged teacher and give the student a height map, so the perceptive policy can lift a foot before contact while a proprioceptive policy adapts only after the step is hit.
\textcite{ross2011dagger} is the reduction that makes that imitation on-policy.
None of those teachers is the worse-informed policy.
Here the PD sees less than the student, and copying it on the approach to a step is the failure mode.

\textcite{rajeswaran2018dapg} keep the task return and use demonstrations as an auxiliary.
\textcite{zhang2021cail} down-weight imitation where the demonstrator is a poor match to the task.
That is the reason the match to the PD is a function of the forward scan rather than a constant label.
\textcite{kumar2021rma} estimate privileged factors from proprioceptive history after contact; the scan is already in the observation, and a latent of that kind is not used.
\textcite{johannink2019residual} add a learned residual on a hand-designed controller.
Residual actions were not used in cycles~1, 2, or~3.
\textcite{rudin2022walk} is the citation for massively parallel PPO on a terrain curriculum, which is the training regime, not a new reward.
\textcite{lee2024wheeled} show wheel driving on flat ground and a legged gait on mapped obstacles for a wheeled-legged robot.
Their teacher still has the cleaner map, and the obstacle gait is a walk rather than this knee tumble.

\section{WaLTER FS and the Reactive PD}
\label{sec:walter}

The teacher is the full-body PD controller.
It commands pitch, roll, drive, and tumble from body state and the commanded velocity.
It does not read \texttt{terrain\_scan}.
The tumble is gated by speed lag.
Until that gate opens, the PD keeps driving the wheels.

On the cycle~0 clips the flat PD holds about \(0.697\,\mathrm{m/s}\) with quiet knees and the gate off, and travels \(5.536\,\mathrm{m}\) in \(8\,\mathrm{s}\).
The command is \(1\,\mathrm{m/s}\).
Hard pyramid stairs are crossed by pitching.
The gate stays off, and the clip ends at \(5.492\,\mathrm{m}\) back on the flat cruise.
Hard inverted stairs open the gate at \(0.909\,\mathrm{m}\), after speed has already fallen, and the cell is not cleared (end distance \(2.929\,\mathrm{m}\)).
Hard box grid opens the gate at \(0.903\,\mathrm{m}\), again after speed has fallen, and ends at \(1.603\,\mathrm{m}\) tipped over (uprightness \(0.664\)).
Knee targets on those latched clips wind to about \(19\,\mathrm{rad}\).

The student sees the same body and command channels plus the height scan.
That scan is the only terrain channel that can move the action before the speed-lag gate would open.
Training terrain for cycles~1 through~3 was the existing \texttt{rough-test} mix.
The mix was not changed.

\section{Training Method}
\label{sec:method}

Cycle~1 trained the teacher environment with PPO.
The objective is the scan-gated match to the PD plus a predictive knee-tumble term, as specified in the method note, together with the existing task rewards.
The match stays on when the forward scan does not show a step, and it turns off when the forward height profile shows one.
Through cycle~2 the step measure was the largest upward jump between consecutive occupied forward bins, so a downward face did not open it.
Cycle~3 uses the largest absolute jump, up or down.
The \(0.04\,\mathrm{m}\) to \(0.08\,\mathrm{m}\) ramp was not changed.
The predictive term pays for knee rate in the direction of the command, and it does not wait for the speed-lag integral.
Through cycle~2 that term was also multiplied by positive body speed, so a stall zeroed it.
Cycle~3 still pays when the forward command is active and the scan gate is open, even if measured speed along the command is at or below zero.
A flat scan keeps the gate shut.
The lag term is still present at its existing weight as a recovery signal.
Training did not use the teacher-only reward.
Residual actions were not implemented.
The scan size and the \texttt{rough-test} mix were not changed.
All three runs are seed 42 only.

The cycle-1 command that finished was
\begin{quote}
\texttt{uv run python scripts/train.py --teacher --task rough-test --num-envs 2048 --max-iterations 500}
\end{quote}
An attempt with 4096 environments failed to allocate \(6366953472\) bytes on the GPU and was discarded.
A two-iteration smoke test with 8 environments exited 0 before that long run.
That checkpoint is \texttt{model\_499.pt} under \texttt{research/logs/training/walter\_fs\_teacher\_rough-test/2026-10-13\_23-32-02/}.

Cycle~2 tested the idea that hard inverted stairs are a downward face, so the upward-only step never opens the predictive term.
The cycle-1 forward-scan clearance falls before the stall, which is less air and therefore higher terrain, not a drop.
A probe of the hard inverted cell reads a \(+0.180\,\mathrm{m}\) riser, match gate 0, and tumble gate 1 from the first frame.
The same probe on hard pyramid stairs reads step \(0.000\,\mathrm{m}\) and tumble gate 0, because that spawn is the top platform and the face ahead goes down.
The gate already opens on the cell that failed, and the cycle-1 policy still moved its knees late.
The gate was left unchanged, and training resumed:
\begin{quote}
\texttt{uv run python scripts/train.py --teacher --task rough-test --num-envs 2048 --max-iterations 500 --resume research/logs/training/walter\_fs\_teacher\_rough-test/2026-10-13\_23-32-02/model\_499.pt}
\end{quote}
The cycle-2 checkpoint is \texttt{model\_998.pt} under \texttt{research/logs/training/walter\_fs\_teacher\_rough-test/2026-10-14\_00-33-54/}.
The file's iteration field is 998, the last of 500 additional updates (iterations 499 through 998).

Cycle~3 tested the cycle-2 inverted-stair frames before editing.
Knee rate there crosses \(1\,\mathrm{rad/s}\) at \(0.832\,\mathrm{m}\) while speed is still \(1.050\,\mathrm{m/s}\), so speed does not go negative before the first spin.
Forward-scan height is already \(0.222\,\mathrm{m}\) at \(0.361\,\mathrm{m}\) of travel, against a flat cruise of \(0.255\,\mathrm{m}\) in that run, so a step is in the scan before the stall.
Speed then reaches \(-0.001\,\mathrm{m/s}\) at \(0.932\,\mathrm{m}\) with knee rate already \(2.257\,\mathrm{rad/s}\).
The old speed factor was zero on that frame.
An import check of the edited gate opened the tumble term on both a \(+0.180\,\mathrm{m}\) rise and a \(-0.180\,\mathrm{m}\) drop, left a flat profile shut, and left a \(0.045\,\mathrm{m}\) slope at match \(0.875\).
Training then resumed:
\begin{quote}
\texttt{uv run python scripts/train.py --teacher --task rough-test --num-envs 2048 --max-iterations 500 --resume research/logs/training/walter\_fs\_teacher\_rough-test/2026-10-14\_00-33-54/model\_998.pt}
\end{quote}
The cycle-3 checkpoint is \texttt{model\_1497.pt} under \texttt{research/logs/training/walter\_fs\_teacher\_rough-test/2026-10-14\_01-31-29/}.
The file's iteration field is 1497, the last of 500 additional updates (iterations 998 through 1497).

\section{One-Cell Protocol}
\label{sec:protocol}

Success is judged on one-cell clips, not on the reward curve.
The evaluator places the robot on one \texttt{rough-test} cell, commands \(1\,\mathrm{m/s}\) forward, records \(8\,\mathrm{s}\), and writes a per-frame table of forward speed, knee rate, uprightness, distance traveled, forward-scan height, student action, PD action, and the PD tumble gate.
Cycle~3 evaluated only the new checkpoint.
The PD numbers are the cycle~0 clips.
Cycles~1 and~2 are the comparison for the same cells.
The old 17-action checkpoint was not retried.

The clips are flat medium, flat hard, hard pyramid stairs, hard inverted stairs, hard box grid, and hard pyramid stairs with the policy height scan replaced by zeros.
Student knee-rate onset is the first frame after \(0.5\,\mathrm{s}\) with knee rate above \(1\,\mathrm{rad/s}\).
The first half-second is a startup transient on flat ground as well.
The PD reference onset is the baseline gate distance.
Flat speed is scored against the \(1\,\mathrm{m/s}\) command.
The PD's \(0.697\,\mathrm{m/s}\) cruise is reported as a gap and is not, by itself, a flat failure.
Hacking signs that fail a cycle even if reward rose are knee spin on flat ground, a high teacher match on stairs with no clearance, spinning in place, falling or bouncing for speed, and a scan ablation that leaves the clip unchanged.
None of those signs is present in the cycle-3 clips.
The goal still fails for the reasons in Section~\ref{sec:results}.

\section{Results}
\label{sec:results}

Cycle~3 resumes the cycle-2 checkpoint.
The iteration-998 line in the new log is \(2.33\), with episode length \(39.6\), because episode lengths were re-randomized at the start of the resume.
By iteration 1048 the mean reward is \(52.12\) and the episode length is \(947.0\).
At iteration 1497 the mean reward is \(51.90\), episode length is \(920.8\), the predictive knee term is \(0.129\), linear-velocity tracking is \(1.509\), and the teacher match is \(0.501\).
Teacher-action MSE at that iteration is \(4.30\).
Terrain level mean is \(2.12\) of a maximum of \(3\).
Iteration 1496 mean reward was \(53.61\), so the last point is inside the late-run band.
Cycle~2 ended at reward \(53.26\).
The curve is cited here and does not decide a pass.
Table~\ref{tab:goal-checks} is the decision.

\begin{table}[ht]
  \centering
  \caption{Cycle~3 goal checks against the PD clips in the baseline review.
    One seed. The goal is not passed.}
  \label{tab:goal-checks}
  \begin{tabular}{ll}
    \toprule
    Check & Status \\
    \midrule
    Flat, medium and hard: knee rate near zero; speed near \(1\,\mathrm{m/s}\) & pass \\
    Hard inverted stairs and hard box grid: clearance at least the PD; earlier onset & fail \\
    Hard pyramid stairs cleared (pitch, not a latched tumble) & pass \\
    Scan ablation on hard pyramid stairs changes knee-rate onset & pass \\
    Two seeds agree & fail \\
    \bottomrule
  \end{tabular}
\end{table}

On flat ground, after the first second, knee-rate means are \(0.016\,\mathrm{rad/s}\) (medium) and \(0.014\,\mathrm{rad/s}\) (hard).
Every frame after \(1\,\mathrm{s}\) is at or below \(0.058\,\mathrm{rad/s}\).
Uprightness stays at \(1.000\).
The gate stays off.
Forward speed stays between \(0.994\) and \(1.070\,\mathrm{m/s}\) on medium and between \(0.995\) and \(1.075\,\mathrm{m/s}\) on hard.
Both means are \(1.027\,\mathrm{m/s}\).
The command error is about \(+0.027\,\mathrm{m/s}\).
The gap versus the PD mean of \(0.697\,\mathrm{m/s}\) is about \(+0.33\,\mathrm{m/s}\).
Cycle~2 means were \(1.093\) and \(1.094\,\mathrm{m/s}\) with knee-rate means \(0.044\,\mathrm{rad/s}\).
Distance at \(8\,\mathrm{s}\) is \(8.190\,\mathrm{m}\) (medium) and \(8.187\,\mathrm{m}\) (hard), against the PD's \(5.536\,\mathrm{m}\).
Frames show level wheel driving, not folded legs.
Head-knee targets near \(0.6\,\mathrm{rad}\) are a held pose. The rate stays near zero.

On hard pyramid stairs the PD gate stays off for all 400 frames.
Knee rate crosses \(1\,\mathrm{rad/s}\) at \(0.680\,\mathrm{m}\), with scan height \(0.431\,\mathrm{m}\) and speed \(0.899\,\mathrm{m/s}\).
Cycle~2 crossed at \(1.235\,\mathrm{m}\).
The PD crosses \(1\,\mathrm{rad/s}\) at \(1.390\,\mathrm{m}\) and also never opens the gate.
Knee commands peak at \(2.019\,\mathrm{rad}\) and return to the flat offset.
Knee rate is back under \(0.1\,\mathrm{rad/s}\) by \(4.379\,\mathrm{m}\).
The student ends at \(7.957\,\mathrm{m}\), speed \(1.006\,\mathrm{m/s}\), knee rate \(0.027\,\mathrm{rad/s}\), uprightness \(1.000\).
The PD ends at \(5.492\,\mathrm{m}\).
The cell is cleared.
The motion is a pitch: the knees bend with the body and release.
Frames go from a level body on the steps, to a nose-down pitch on the stair face, to a level body on the gray surround.

On hard inverted stairs, knee rate crosses \(1\,\mathrm{rad/s}\) at \(0.718\,\mathrm{m}\) with speed still \(1.008\,\mathrm{m/s}\) and uprightness \(1.000\).
Speed falls below \(0.5\,\mathrm{m/s}\) later, at \(0.765\,\mathrm{m}\), while knee rate is already \(2.511\,\mathrm{rad/s}\).
The PD opens its gate at \(0.909\,\mathrm{m}\) after speed has already fallen.
Cycle~2 crossed at \(0.832\,\mathrm{m}\) with speed \(1.050\,\mathrm{m/s}\).
Onset is earlier than the PD.
End distance is \(6.820\,\mathrm{m}\) against the PD's \(2.929\,\mathrm{m}\), with knee rate \(1.401\,\mathrm{rad/s}\) and uprightness \(0.826\).
Peak knee command after \(0.5\,\mathrm{s}\) is \(10.548\,\mathrm{rad}\), against cycle~2's \(8.022\,\mathrm{rad}\) and the PD's windup to about \(19\,\mathrm{rad}\).
Frames at \(5\,\mathrm{s}\) and \(7\,\mathrm{s}\) are still on the red stair face, short of the gray surround.
The cell is not cleared.
The knees do not return to the flat offset, so this is a tumble attempt rather than a pitch that releases.

On the hard box grid, knee rate crosses \(1\,\mathrm{rad/s}\) at \(0.450\,\mathrm{m}\), speed \(0.507\,\mathrm{m/s}\), uprightness \(0.969\), scan height \(0.514\,\mathrm{m}\), gate off.
The PD opens its gate at \(0.903\,\mathrm{m}\) with speed already \(0.256\,\mathrm{m/s}\).
Cycle~2 crossed at \(0.431\,\mathrm{m}\) and then left the cell.
The student ends at \(5.320\,\mathrm{m}\), speed \(-0.186\,\mathrm{m/s}\), knee rate \(2.477\,\mathrm{rad/s}\), uprightness \(0.997\), gate on.
The PD ends at \(1.603\,\mathrm{m}\), uprightness \(0.664\).
Peak knee command is \(9.603\,\mathrm{rad}\), and the end commands are still several radians.
Frames at \(5\,\mathrm{s}\) and \(7\,\mathrm{s}\) are still in the green grid, legs folded, short of the gray surround.
Distance is past the PD.
The cell is not cleared, so the combined check fails.
Cycle~2 had ended upright on the gray surround at \(6.634\,\mathrm{m}\).

Zeroing the height scan on hard pyramid stairs removes the knee-rate crossing.
With the scan it is at \(0.680\,\mathrm{m}\).
Without the scan, knee rate after \(0.5\,\mathrm{s}\) stays at or below \(0.817\,\mathrm{rad/s}\).
Final distance changes from \(7.957\,\mathrm{m}\) to \(12.800\,\mathrm{m}\), and final speed from \(1.006\,\mathrm{m/s}\) to \(1.613\,\mathrm{m/s}\).
The uprightness drop is at \(1.125\,\mathrm{m}\) with the scan and at \(1.273\,\mathrm{m}\) without it.
Both PD gates stay off.
Cycle~2 moved the crossing from \(1.235\,\mathrm{m}\) to \(4.504\,\mathrm{m}\), and distance from \(8.345\,\mathrm{m}\) to \(10.399\,\mathrm{m}\).
The scan changes the clip.
The with-scan frame at \(1\,\mathrm{s}\) is still level; the ablated frame at \(1\,\mathrm{s}\) is already pitched.

\section{Failure Cases}
\label{sec:failures}

The flat policy is upright and the knees are quiet, and the speed sits near the command.
The mean speed \(1.027\,\mathrm{m/s}\) is about \(0.027\,\mathrm{m/s}\) above \(1\,\mathrm{m/s}\) and about \(0.33\,\mathrm{m/s}\) above the PD mean of \(0.697\,\mathrm{m/s}\).
The PD gap is the controller's undershoot.
It is not a flat failure.
Knee rate after \(1\,\mathrm{s}\) stays at or below \(0.058\,\mathrm{rad/s}\).

Hard pyramid stairs are cleared by pitching.
The PD gate never opens.
Knee commands peak at \(2.019\,\mathrm{rad}\) and return.
Knee-rate onset at \(0.680\,\mathrm{m}\) is earlier than cycle~2's \(1.235\,\mathrm{m}\), and the motion still releases.
Hard inverted stairs start the knees before the stall, at \(0.718\,\mathrm{m}\) against the PD gate at \(0.909\,\mathrm{m}\), and still do not leave the stair face.
Distance \(6.820\,\mathrm{m}\) is past the PD's \(2.929\,\mathrm{m}\).
The \(7\,\mathrm{s}\) frame is still on the red stairs, and uprightness at \(8\,\mathrm{s}\) is \(0.826\).
The hard box grid also starts earlier, at \(0.450\,\mathrm{m}\) against the PD gate at \(0.903\,\mathrm{m}\), and does not leave the grid.
The \(7\,\mathrm{s}\) frame is still among the green cubes.
End speed is \(-0.186\,\mathrm{m/s}\) and the PD gate is still on.
Cycle~2 had cleared that cell.
Both obstacle cells therefore fail the combined check, even though both onsets are earlier than the PD.

The scan ablation changes onset and distance, so the policy is not ignoring the scan on hard pyramid stairs.
With the scan the crossing is at \(0.680\,\mathrm{m}\).
Without it, knee rate after \(0.5\,\mathrm{s}\) never reaches \(1\,\mathrm{rad/s}\).
Both PD gates stay off.
The ablated body is already pitched at \(1\,\mathrm{s}\), while the scan clip at \(1\,\mathrm{s}\) is still level.

No clip shows the hacking signs named in the protocol.
Flat knees do not keep spinning, the pyramid clip is not a frozen copy of a late PD, distance traveled is several meters on every cell, and the ablated clip is not a copy of the scan clip.
The cycle fails the goal checks anyway, because hard inverted stairs and the hard box grid are not cleared and only seed 42 was run.
Two-seed agreement is not met and cannot be claimed.
The study is unfinished.

\section{Conclusion}
\label{sec:conclusion}

This is cycle~3 of an unfinished study.
The PD baseline remains the reactive reference: flat driving at about \(0.70\,\mathrm{m/s}\) with quiet knees, pyramid stairs cleared by pitching with the gate off, and inverted stairs and the box grid opened only after speed has already fallen.
The cycle-2 inverted clip already spun the knees at positive speed, at \(0.832\,\mathrm{m}\), and then stalled at the stair with knee rate \(2.257\,\mathrm{rad/s}\) while forward speed was \(-0.001\,\mathrm{m/s}\).
Cycle~3 stopped zeroing the predictive term on that stall and counted drops in the step measure.
The seed-42 student was trained for 500 further iterations at 2048 environments, from \texttt{model\_998.pt} to \texttt{model\_1497.pt}, and mean reward at iteration 1497 is \(51.90\).
Flat knee rate is near zero and flat speed is near the \(1\,\mathrm{m/s}\) command, about \(0.33\,\mathrm{m/s}\) above the PD.
Hard pyramid stairs are cleared by pitching.
Hard inverted stairs start the knee-rate rise at \(0.718\,\mathrm{m}\) against the PD gate at \(0.909\,\mathrm{m}\) and do not leave the stairs.
Hard box grid starts at \(0.450\,\mathrm{m}\) against \(0.903\,\mathrm{m}\) and also does not leave the cell, which cycle~2 had cleared.
A scan ablation moves hard-pyramid knee-rate onset from \(0.680\,\mathrm{m}\) to never.
A second seed was not trained.
The paper goal is not passed.

\printbibliography

\end{document}
